\section{Java Library}
\label{sec:JavaLibrary}

The Java implementation is a Maven reactor. Application developers normally
interact with the runtime and state-machine modules rather than constructing
\texttt{RaftNode} directly.

\subsection{Application integration path}

A Java application should start from:

\begin{itemize}
\item \texttt{SnapshotStateMachine} or \texttt{QueryableStateMachine} in
\filepath{raft-state-machine},
\item \texttt{BasicAdapter} and application module hooks in \filepath{raft-runtime},
\item a transport factory, normally Netty through \filepath{raft-transport-netty},
\item storage from \filepath{raft-storage}.
\end{itemize}

The normal shape is:

\begin{enumerate}
\item Define application command and query payloads.
\item Implement the state machine.
\item Create an adapter or module that supplies that state machine to the runtime.
\item Configure peers, local identity, storage and transport.
\item Let the runtime expose client command/query and operational endpoints.
\end{enumerate}

\subsection{State machine interfaces}

\literalbreak{SnapshotStateMachine} is the minimal contract. It receives committed
commands and can produce/restore application snapshots. \literalbreak{QueryableStateMachine}
adds local query support. The result-returning state-machine variant is useful
when a committed command should return a typed application result, for example
CAS.

State-machine implementations should be thread-conscious. Raft calls are
synchronized at the node level where necessary, but application state should
still avoid exposing mutable internal collections directly.

\subsection{Runtime modules}

\texttt{RaftApplicationModule} separates application packaging from the shared
runtime. It lets an application declare supported runtime modes, CLI commands
and an adapter factory. This is where reference-data or other domain packages
plug into the generic executable.

\texttt{RaftApplicationFactory} creates a runnable adapter for a local peer. It
receives:

\begin{itemize}
\item local peer identity,
\item initially known peers,
\item optional join seed,
\item runtime configuration,
\item timeout settings.
\end{itemize}

An application should prefer these runtime hooks over direct construction of
transport and node internals.

\subsection{Policy hooks}

The Java runtime exposes:

\begin{itemize}
\item \texttt{ClientWriteAdmissionPolicy},
\item \texttt{ClientCommandAuthenticator},
\item \texttt{ClientCommandAuthorizer}.
\end{itemize}

Use these for domain and security decisions before a command enters Raft. Keep
the policy decisions explicit and testable. Do not hide domain rejection inside a
state-machine no-op unless that no-op is itself the deterministic domain
behavior.

\subsection{Storage}

The Java storage layer separates persistent term/vote state from the log store.
Production deployments should use durable stores. Tests may use the in-memory
stores, but those stores are not a recovery mechanism.

Cluster configuration is applied from internal replicated entries and is also
carried in snapshot metadata. Application code should treat configuration as
Raft-owned state.

\subsection{Javadocs and Maven publishing}

The reactor is configured to attach source and Javadoc jars. The public API
should therefore keep class and method Javadocs useful for application
developers. In particular, boundary interfaces should document whether a method
is application-owned or Raft-owned.

The usual publishing readiness checks are:

\begin{itemize}
\item \shellcmd{mvn -q test}
\item \shellcmd{mvn -q -DskipTests package}
\item verify source and Javadoc jars under module \filepath{target} directories
\item ensure publishing metadata and signing are configured before Maven Central upload
\end{itemize}
